% Figura (Problema 1.3): hilo infinito (geometría) e hilo finito frente a la aproximación de hilo infinito.
\begin{figure}[ht]
\centering
\pgfmathdeclarefunction{Erfin}{2}{%
  \begingroup
  \pgfmathparse{(1/(#1))*(((#2)+0.5)/sqrt((#1)*(#1)+((#2)+0.5)*((#2)+0.5))-((#2)-0.5)/sqrt((#1)*(#1)+((#2)-0.5)*((#2)-0.5)))}%
  \pgfmathsmuggle\pgfmathresult\endgroup}
\pgfmathdeclarefunction{Ezfin}{2}{%
  \begingroup
  \pgfmathparse{-(1/sqrt((#1)*(#1)+((#2)+0.5)*((#2)+0.5))-1/sqrt((#1)*(#1)+((#2)-0.5)*((#2)-0.5)))}%
  \pgfmathsmuggle\pgfmathresult\endgroup}
\pgfmathdeclarefunction{errfin}{2}{%
  \begingroup
  \pgfmathparse{100*abs(2/(#1)-sqrt(Erfin(#1,#2)*Erfin(#1,#2)+Ezfin(#1,#2)*Ezfin(#1,#2)))/sqrt(Erfin(#1,#2)*Erfin(#1,#2)+Ezfin(#1,#2)*Ezfin(#1,#2))}%
  \pgfmathsmuggle\pgfmathresult\endgroup}
\begin{minipage}[t]{0.36\textwidth}\centering
\begin{tikzpicture}[baseline=(current bounding box.north), x=1cm, y=1cm]
  \draw[eje] (-0.2,0) -- (3.4,0) node[below] {\small $x$};
  \draw[line width=2.2pt, azulprof] (0,-2.0) -- (0,2.3);
  \draw[eje] (0,2.3) -- (0,2.8) node[above] {\small $z$};
  \fill[azulprof] (0,0) circle (1.4pt);
  \coordinate (P) at (2.3,0);
  \coordinate (A) at (0,1.2); \coordinate (B) at (0,-1.2);
  \fill[azulacero] (A) circle (2pt) node[left, font=\small] {$\dif q'$};
  \fill[azulacero] (B) circle (2pt) node[left, font=\small] {$\dif q'$};
  \draw[->, azulclaro!70!azulacero, line width=0.7pt] (A) -- (P) node[pos=0.45, above right, font=\small, text=azulprof] {$\vc{R}$};
  \draw[->, azulclaro!70!azulacero, line width=0.7pt] (B) -- (P);
  \draw[vectrojo] (P) -- ++(0.55,0.28) node[above right, font=\scriptsize, text=acento] {$\dif\vE_{1}$};
  \draw[vectrojo] (P) -- ++(0.55,-0.28) node[below right, font=\scriptsize, text=acento] {$\dif\vE_{2}$};
  \draw[vect, azulprof, line width=1.2pt] (P) -- ++(1.1,0) node[right, font=\small] {$\vE$};
  \draw[grispizarra, densely dotted] (0,0) -- (P) node[midway, below, font=\small, text=azulprof] {$r$};
  \node[font=\scriptsize, text=grispizarra, anchor=west] at (-1.5,1.2) {$z'$};
  \node[font=\scriptsize, text=grispizarra, anchor=west] at (-1.5,-1.2) {$-z'$};
\end{tikzpicture}\\[2pt]{\small\scshape (a)}
\end{minipage}\hfill
\begin{minipage}[t]{0.62\textwidth}\centering
\begin{tikzpicture}[baseline=(current bounding box.north), x=7.6cm, y=2.4cm]
  \draw[grispizarra, line width=0.5pt] (-0.55,-1.05) rectangle (0.55,1.05);
  \foreach \xx in {-0.45,-0.35,...,0.46}{\foreach \zz in {-0.9,-0.8,...,0.91}{%
    \pgfmathsetmacro{\ra}{abs(\xx)}%
    \pgfmathsetmacro{\er}{Erfin(\ra,\zz)}\pgfmathsetmacro{\ez}{Ezfin(\ra,\zz)}%
    \pgfmathsetmacro{\em}{sqrt(\er*\er+\ez*\ez)}%
    \pgfmathsetmacro{\ft}{min(1,(\em^0.4)/3.2)}%
    \mapcolor{\ft}%
    \pgfmathsetmacro{\ux}{sign(\xx)*\er/\em}\pgfmathsetmacro{\uz}{\ez/\em}%
    \pgfmathsetmacro{\ll}{0.03*(0.55+\ft)}%
    \draw[-{Stealth[length=3pt,width=3pt]}, line width=0.5pt, mapc] (\xx-\ux*\ll*0.9,\zz-\uz*\ll*0.9/3.42) -- ++({\ux*\ll*1.8},{\uz*\ll*1.8/3.42});}}
  \draw[line width=2.2pt, azulprof] (0,-0.5) -- (0,0.5);
  \foreach \zz in {0,0.1,0.2,0.3}{\draw[azulprof, densely dashed, line width=0.5pt] (-0.5,\zz) -- (0.5,\zz);
     \node[font=\tiny, text=azulprof, fill=white, inner sep=0.6pt, anchor=west] at (0.505,\zz) {\pgfmathprintnumber[fixed,precision=1,use comma]{\zz}};}
  \foreach \t in {-0.4,-0.2,0,0.2,0.4}{\node[font=\tiny, anchor=north] at (\t,-1.05) {\pgfmathprintnumber[fixed,precision=1,use comma]{\t}};}
  \node[font=\scriptsize, anchor=north] at (0,-1.25) {$x/L$};
  \foreach \t in {-1,-0.5,0,0.5,1}{\node[font=\tiny, anchor=east] at (-0.555,\t) {\pgfmathprintnumber[fixed,precision=1,use comma]{\t}};}
  \node[font=\scriptsize, rotate=90, anchor=south] at (-0.66,0) {$z/L$};
\end{tikzpicture}\\[2pt]{\small\scshape (b)}
\end{minipage}\\[6pt]
\pgfplotsset{errax/.style={tema, width=0.265\textwidth, height=3.3cm, xmin=0, xmax=0.5, ymin=0, ymax=100,
   xtick={0,0.25,0.5}, ytick={0,50,100}, xlabel={$r/L$}, domain=0.02:0.5, samples=60, grid=major}}
\begin{tikzpicture}
  \begin{axis}[errax, name=a, ylabel={error (\%)}, title={\scriptsize $z/L=0{,}0$}]
    \addplot[azulprof, line width=1.0pt] {errfin(x,0.0)};
  \end{axis}
  \begin{axis}[errax, at={(a.east)}, anchor=west, xshift=0.35cm, name=b, yticklabels={}, title={\scriptsize $z/L=0{,}1$}]
    \addplot[azulprof, line width=1.0pt] {errfin(x,0.1)};
  \end{axis}
  \begin{axis}[errax, at={(b.east)}, anchor=west, xshift=0.35cm, name=c, yticklabels={}, title={\scriptsize $z/L=0{,}2$}]
    \addplot[azulprof, line width=1.0pt] {errfin(x,0.2)};
  \end{axis}
  \begin{axis}[errax, at={(c.east)}, anchor=west, xshift=0.35cm, yticklabels={}, title={\scriptsize $z/L=0{,}3$}]
    \addplot[azulprof, line width=1.0pt] {errfin(x,0.3)};
  \end{axis}
\end{tikzpicture}\\[2pt]{\small\scshape (c)}
\caption{Problema 1.3. (a)~Hilo infinito: los elementos $\dif q'$ en $z'$ y $-z'$ crean campos cuyas componentes $z$ se cancelan y
solo sobrevive la radial. (b)~Campo de un hilo \emph{finito} de longitud $L$ (trazo grueso) en el plano $(x,z)$, en unidades
de $\lambda/4\pi\eo L$ (el color y la longitud de las flechas codifican $\norm{\vE}$, de forma comprimida). (c)~Error relativo
entre el módulo del campo exacto del hilo finito y el de la aproximación de hilo infinito, $\lambda/2\pi\eo r$, a lo largo de las rectas horizontales
$z/L=0;\,0{,}1;\,0{,}2;\,0{,}3$ marcadas a trazos en (b): crece al alejarnos, pero la aproximación es muy buena cerca del hilo y lejos de sus extremos.}
\label{fig:sem-hilo}
\end{figure}
